\section{Criticità}

Durante il progetto, come già descritto in precedenza, si è passati da una suddivisione paritaria per tutte le fasi a una ripartizione basata sulle competenze pregresse. Ciò ha comportato un aumento del carico di lavoro per \textit{Marco Piovesan} e \textit{Alessio Pesaresi} nella prima parte, a fronte di un impegno minore per \textit{Michele Porcellato} nella fase di \texttt{D1}; quest'ultimo ha tuttavia sostenuto un carico decisamente superiore nella fase di \texttt{D2} (scrittura del codice), in qualità di sviluppatore principale dell'applicazione.

Inoltre, mentre durante i periodi di lezione si è riusciti a mantenere un ritmo operativo medio-basso, l'avanzamento dei lavori ha subito due forti battute d'arresto in corrispondenza delle sessioni d'esame. Una prima flessione si è verificata durante la pausa natalizia e la sessione di gennaio/febbraio a causa della preparazione dei relativi esami, richiedendo un successivo sforzo di recupero. Una dinamica del tutto analoga si è poi ripresentata nella sessione estiva tra giugno e luglio: il sovrapporsi delle prove accademiche ha rallentato nuovamente lo sviluppo, rendendo irrealizzabile l'obiettivo originario di consegnare il progetto a luglio e costringendo il gruppo a posticipare la consegna finale alla sessione di settembre.

Un'ulteriore criticità rilevante è emersa durante la stesura del \texttt{D2} e lo sviluppo del software: è stata infatti riscontrata una marcata discrepanza tra le specifiche delineate nel documento \texttt{D1} e le reali scelte architetturali emerse in fase implementativa. Tale disallineamento ha imposto una revisione sostanziale del documento \texttt{D1}, rallentando temporaneamente le attività ma garantendo una documentazione finale coerente e fedele all'architettura realizzata.

Sebbene i componenti del gruppo non abbiano contribuito in misura identica a ciascun \textit{deliverable}, l'impegno complessivo è risultato ben bilanciato e tutti i membri hanno partecipato attivamente al successo del progetto. Chi non è stato direttamente coinvolto nella stesura di uno specifico documento ha comunque contribuito alla sua fase di controllo e correzione: ciascun membro ha revisionato almeno una volta ogni \textit{deliverable} alla cui redazione non aveva preso parte attiva, assicurando così l'accuratezza e l'omogeneità del risultato finale.